%%%%%%%%%%%%%%%%%%%%%%%%%%%%%%%%%%%%%%%%%%%%%%%%%%%%%%%%%%%%%%%%%%%%%%%%%%%%%%%%%
%	TASK 3 - Evaluation Methodology and Analysis
%   Maps to brief Tasks 3.1, 3.2, 3.3.
%%%%%%%%%%%%%%%%%%%%%%%%%%%%%%%%%%%%%%%%%%%%%%%%%%%%%%%%%%%%%%%%%%%%%%%%%%%%%%%%%
\lhead{Task 3. \emph{Evaluation Methodology and Analysis}}
\section{Task 3: Evaluation Methodology and Analysis}
\label{Task3}

\subsection{Parameter Sensitivity Analysis (Task 3.1)}
\label{T3:Sensitivity}
TODO: Investigate how GA parameters influence performance. Report configurations that deliver:
\begin{itemize}
    \item robust performance;
    \item reasonable computational cost;
    \item consistency across multiple runs.
\end{itemize}

\subsection{Performance Evaluation (Task 3.2)}
\label{T3:Performance}
TODO: Evaluate the GA on:
\begin{itemize}
    \item solution quality (tour length);
    \item convergence rate;
    \item stability across trials;
    \item computational efficiency.
\end{itemize}

\subsection{Critical Analysis of Repair Mechanism (Task 3.3)}
\label{T3:RepairCritique}
This section evaluates the repair mechanism using the controlled comparison
from TICKET-17. Both experimental arms used the same GA settings on kroA100:
population size 100, 500 generations, tournament selection, naive single-point
crossover, swap mutation, and 30 independent seeds. The only experimental
difference was whether duplicate/missing-city repair was applied after
crossover and mutation.

\begin{table}[h]
    \centering
    \caption{Sanity check showing that naive crossover creates infeasibility and that repair restores a valid permutation.}
    \label{tab:repair-sanity}
    \begin{tabular}{l r}
        \hline
        Observation from TICKET-17 sanity check & Value \\
        \hline
        Parent A valid permutation & True \\
        Parent B valid permutation & True \\
        Naive-crossover child valid permutation & False \\
        Duplicate positions in child & 23 \\
        Missing cities in child & 23 \\
        Repaired child valid permutation & True \\
        \hline
    \end{tabular}
\end{table}

\begin{table}[h]
    \centering
    \caption{Final-generation summary over 30 paired runs per strategy. Lower best fitness is numerically smaller, but the no-repair arm is infeasible under the TSP constraint.}
    \label{tab:repair-summary}
    \begin{tabular}{l r r r r r}
        \hline
        Strategy & Mean best & Std. & Min. best & Gap to optimum & Final diversity \\
        \hline
        No repair & 38129.27 & 4135.41 & 28805.71 & 79.2\% & 0.0100 \\
        Repair & 49777.84 & 3186.86 & 43890.51 & 133.9\% & 0.0128 \\
        \hline
    \end{tabular}
\end{table}

\begin{table}[h]
    \centering
    \caption{Mean best-fitness and diversity checkpoints from TICKET-17 convergence logs. Values are averaged over 30 seeds per strategy.}
    \label{tab:repair-convergence}
    \begin{tabular}{l r r r}
        \hline
        Strategy & Generation & Mean best fitness & Mean diversity \\
        \hline
        No repair & 0 & 151041.61 & 0.9900 \\
        No repair & 100 & 61144.48 & 0.0360 \\
        No repair & 250 & 43844.78 & 0.0165 \\
        No repair & 500 & 38129.27 & 0.0100 \\
        Repair & 0 & 151041.61 & 0.9900 \\
        Repair & 100 & 77552.75 & 0.0413 \\
        Repair & 250 & 59334.24 & 0.0247 \\
        Repair & 500 & 49777.84 & 0.0128 \\
        \hline
    \end{tabular}
\end{table}

\begin{table}[h]
    \centering
    \caption{Paired Wilcoxon signed-rank test on final best fitness from the 30 matched seeds in TICKET-17.}
    \label{tab:repair-wilcoxon}
    \begin{tabular}{l r}
        \hline
        Test quantity & Value \\
        \hline
        Statistic & 0.00 \\
        $p$-value & $1.86 \times 10^{-9}$ \\
        Significance level & 0.05 \\
        Decision & Reject equal medians \\
        Mean paired difference, repair $-$ no repair & 11648.57 \\
        \hline
    \end{tabular}
\end{table}

The repair mechanism is effective at enforcing feasibility in the scenario for
which it was designed. The TICKET-17 sanity check shows that naive crossover
changed two valid parents into an invalid child with 23 duplicate positions and
23 missing cities, while the repair step converted that same child back into a
valid permutation (\tref{tab:repair-sanity}). This establishes that the repair
operator is not cosmetic in this experiment: without it, the GA evaluates
chromosomes that do not represent Hamiltonian tours (\tref{tab:repair-sanity}).

The final-fitness results must be interpreted through feasibility, not raw
numeric distance alone. The no-repair arm obtained a lower mean best fitness
than the repair arm, 38129.27 versus 49777.84 (\tref{tab:repair-summary}), and
the paired Wilcoxon signed-rank test reported a test statistic of 0.00 with
$p = 1.86 \times 10^{-9}$ (\tref{tab:repair-wilcoxon}). This does not show that
no-repair is a better TSP solver. Instead, it confirms the infeasibility
pathology: when duplicate cities are allowed, the chromosome can skip some
cities and still receive a short distance from the fitness function
(\tref{tab:repair-sanity}, \tref{tab:repair-summary}). The repair arm reports
higher distances because it evaluates genuine 100-city tours after constraint
restoration (\tref{tab:repair-summary}).

The main trade-off is between strict feasibility and exploratory freedom. Repair
removes invalid chromosomes before evaluation, which protects selection pressure
from artificially short infeasible tours (\tref{tab:repair-summary}). However,
repair also edits offspring after crossover and mutation, so the relevant
empirical question is whether that editing caused faster diversity collapse in
the observed runs. In the TICKET-17 logs, both strategies lost most of their
initial Hamming diversity by generation 500, but the repair arm retained
slightly higher final diversity than no-repair (0.0128 versus 0.0100) and
stayed higher at the generation-100 and generation-250 checkpoints
(\tref{tab:repair-convergence}). This suggests that, for the random-insertion
repair used here, feasibility enforcement did not cause faster diversity
collapse than the no-repair baseline in this run (\tref{tab:repair-convergence}).

Compared with alternative constraint-handling strategies, repair is most useful
when the chosen operators can create invalid permutations. A feasibility-
preserving crossover such as PMX, OX, or CX avoids this specific duplicate/missing
city failure by construction, so the repair step would usually become a safeguard
rather than the main constraint handler. By contrast, the no-repair baseline in
TICKET-17 is equivalent to using no effective penalty for infeasibility: invalid
chromosomes survive and are selected because their raw distances are shorter
(\tref{tab:repair-summary}). A penalty-function alternative would need to assign
enough cost to duplicate and missing cities to prevent this selection-pressure
distortion; otherwise it would reproduce the same misleading advantage observed
in the no-repair arm (\tref{tab:repair-sanity}, \tref{tab:repair-summary}).

Overall, TICKET-17 supports the repair mechanism as a necessary constraint-
handling component for naive crossover experiments. The statistical test shows
that the two strategies behave differently, the sanity check shows why the
difference occurs, and the convergence checkpoints show that repair preserves
validity while maintaining slightly more final diversity than the invalid
baseline in this experiment (\tref{tab:repair-sanity},
\tref{tab:repair-summary}, \tref{tab:repair-convergence},
\tref{tab:repair-wilcoxon}). The repair arm is therefore the meaningful basis
for TSP performance discussion, while the no-repair arm is best interpreted as
evidence of how infeasibility corrupts fitness-based selection.
